\documentclass[11pt,a4paper]{article}
\usepackage[margin=26mm,headheight=14pt]{geometry}
\usepackage{fontspec}\setmainfont{Noto Serif}\setsansfont{Noto Sans}
\usepackage{amsmath,amssymb,booktabs,tabularx,enumitem,xcolor,microtype,fancyhdr,hyperref}
\hypersetup{colorlinks=true,linkcolor=blue!45!black,urlcolor=blue!45!black,pdftitle={Independent mathematical audit of the 3D contour-generator bound}}
\definecolor{darkblue}{HTML}{1B5271}
\definecolor{darkred}{HTML}{A33636}
\definecolor{darkgreen}{HTML}{267446}
\newcommand{\pass}{\textcolor{darkgreen}{\textsf{VERIFIED}}}
\newcommand{\flag}{\textcolor{darkred}{\textsf{FLAG}}}
\newcommand{\qual}{\textcolor{darkblue}{\textsf{QUALIFIED}}}
\pagestyle{fancy}\fancyhf{}\lhead{\small\sffamily INDEPENDENT MATHEMATICAL AUDIT}\rhead{\small\sffamily 4 NOVEMBER 2026}\cfoot{\thepage}
\emergencystretch=3em
\setlength{\parindent}{0pt}\setlength{\parskip}{6pt}
\begin{document}
{\sffamily\Large\bfseries Referee Verification and Editorial Flags}\\[5pt]
{\large On the 3D contour-generator Willmore-type estimate}\\[3pt]
\emph{Source note attributed to İmamverdi Bədəlov. This audit is not a certification of novelty or a machine-checked proof.}
\smallskip\hrule\bigskip
\section*{Decision}
\textbf{Main mathematical statement:} \pass.\\
\textbf{Two-page proof as originally submitted:} \flag{} --- incomplete; repaired in the accompanying paper.\\
\textbf{Sharpness and equality:} \pass{} for the stated smooth, compact, boundaryless setting.\\
\textbf{Projected-length identity:} \pass{} as a mathematical derivation; original attribution and novelty unverified.\\
\textbf{Novelty / journal submission:} \qual{} --- priority research and human authorship review needed.

The audited upper bound is
\[
\boxed{\quad\mathcal L(\Sigma)^2\leq\frac{A(\Sigma)}2\bigl(\mathcal W(\Sigma)-\pi\chi(\Sigma)\bigr),\quad
\mathcal W=\int_\Sigma H^2\,dA,\quad H=\frac{\kappa_1+\kappa_2}{2}.\quad}
\]
The coefficient is exactly attained by a round sphere and thus is optimal. Its proof is valid after replacing the source note's summary of coarea and genericity by an incidence-manifold argument.

\section*{Proof ledger: what was actually checked}
\renewcommand{\arraystretch}{1.25}
\begin{tabularx}{\linewidth}{@{}p{3.3cm}p{2cm}X@{}}\toprule
Claim & Status & Reason / repair \\\midrule
Almost-everywhere regularity & \pass & Define $M=\{(p,u):N(p)\cdot u=0\}$, a smooth compact 3-manifold; apply Sard to $M\to\mathbb S^2$.\\
Exact average spatial rim formula & \pass & Apply coarea to both incidence projections and compute their $2$-Jacobians explicitly.\\
Ellipse perimeter estimate & \pass & Jensen yields $P\leq\pi\sqrt{2(\kappa_1^2+\kappa_2^2)}$ with the required coefficient.\\
Energy identity & \pass & Use $\kappa_1^2+\kappa_2^2=4H^2-2K$ and $\int K=2\pi\chi$.\\
Optimality & \pass & Sphere radius $r$: $\mathcal L=2\pi r$, $A=4\pi r^2$, $\mathcal W=4\pi$, $\chi=2$.\\
Equality characterization & \pass & Jensen equality forces equal squared principal curvatures; Cauchy--Schwarz forces constant $|B|$. Nontrivial compact minimal components cannot occur.\\
Projected contour extension & \pass & The new coarea weight is $|\langle Bu,u\rangle|$, the absolute normal curvature, not $|Bu|$.\\
\bottomrule
\end{tabularx}

\section*{Editorial and mathematical flags}
\begin{enumerate}[leftmargin=*,itemsep=4pt]
\item \textbf{F1 --- Incorrect submersion wording (repaired).} The source calls $p\mapsto N(p)\cdot u$ a submersion; a real-valued smooth function on a compact surface always has critical points. Only the zero level needs to be regular for almost all views.
\item \textbf{F2 --- Unjustified transversality/stratification claim (repaired).} The source appeals to a codimension-one set without constructing the relevant map or proving the stratification. Sard's theorem on the incidence manifold is sufficient.
\item \textbf{F3 --- Missing coarea Jacobian calculation (repaired).} On $M$ the relevant Jacobians are $J_2\pi_p=(1+|Bu|^2)^{-1/2}$ and $J_2\pi_u=|Bu|(1+|Bu|^2)^{-1/2}$.
\item \textbf{F4 --- Degenerate ellipse ambiguity (repaired).} The ellipse perimeter is interpreted as the length of its parameterized image, counting a collapsed segment twice.
\item \textbf{F5 --- Visibility terminology (qualified).} Spatial rim length, projected rim length with multiplicity, and visible image silhouette length are not equal in general. The corrected paper explicitly distinguishes them.
\item \textbf{F6 --- Literature and priority (open).} Existing work defines and studies contour generators and apparent contours. A targeted search did not verify a previous statement of these exact averaged formulas, but absence of a search hit is not evidence of novelty.
\item \textbf{F7 --- Added results and author attribution (open).} Equality classification, lower comparison, and projected identity were developed during this review. They should not be silently presented as content of the original note or submitted without authorial review.
\item \textbf{F8 --- Hypothesis scope (qualified).} The proof covers compact smooth oriented surfaces without boundary in $\mathbb R^3$. A surface with boundary or an unbounded surface would require separate hypotheses or boundary/integrability provisions.
\end{enumerate}

\section*{Independent numerical falsification attempts}
For a ring torus $(R,r)$, $R>r$, we computed both (i) the principal-curvature integrals and (ii) the lengths of directly traced rim curves for generic viewing directions followed by a directional average. Fixed-grid quadrature initially struggled near horizontal views; adaptive quadrature resolves the narrow peaks. The new projected-length formula was tested by the same two independent methods.
\begin{center}\small
\begin{tabular}{@{}lrrr@{}}\toprule
Torus $(R,r)$ & $\mathcal L$ & $\mathcal L_{\mathrm{pr}}$ & bound $\sqrt{A\mathcal W/2}$\\\midrule
$(2,1)$ & 28.23365907 & 23.19516093 & 29.99711051\\
$(\sqrt{2},1)$ & 21.60993419 & 18.17206067 & 23.47400755\\
$(3,0.5)$ & 38.46296443 & 30.57402222 & 42.16912718\\\bottomrule
\end{tabular}
\end{center}
The independently integrated direct rim values differed from the first column by less than $2\times10^{-6}$; projected values differed by less than $10^{-5}$ at the chosen resolution. All are below the proposed upper bound. Sphere tests at radii $0.5$, $1$, and $3$ give exact equality analytically. See the accompanying \texttt{numerical\_verification.py} and its printed results. These tests are supplementary: a finite selection of surfaces cannot establish a theorem.

\section*{Possible strengthening and practical interpretation}
The full rim length obeys the two-sided bound
\[
\frac1\pi\int_\Sigma |B|\,dA\leq\mathcal L(\Sigma)\leq
\frac1{2\sqrt2}\int_\Sigma |B|\,dA.
\]
For projected \emph{parameterized} contours, the independent identity is
\[
\mathcal L_{\mathrm{pr}}(\Sigma)=\frac1{4\pi}\int_\Sigma\int_0^{2\pi}
|\kappa_1\cos^2\theta+\kappa_2\sin^2\theta|\,d\theta\,dA.
\]
It yields $\frac12\int|H|\,dA\leq\mathcal L_{\mathrm{pr}}\leq\mathcal L$; for a convex ovaloid the left inequality becomes an equality. These observations are useful when the motivating application concerns two-dimensional projected contours rather than spatial rim curves.

\section*{What remains unverified}
\textbf{Novelty:} no exhaustive review of the classical integral-geometry or computer-vision literature was performed; no claim of original priority is endorsed. \textbf{Publication:} the manuscript has not been externally peer reviewed. \textbf{Formal verification:} the proofs were audited mathematically and with reproducible numerical examples, but not translated into a proof assistant such as Lean or Coq. \textbf{Applications:} no image-processing benchmarks or visibility-aware formulas have been established.

\section*{Sources used for background checks}
\begin{itemize}[leftmargin=*,itemsep=0pt]
\item H. Federer, \emph{Geometric Measure Theory}, Springer, 1969 (coarea theorem).
\item M. P. do Carmo, \emph{Differential Geometry of Curves and Surfaces}, Prentice--Hall, 1976.
\item S. Plantinga and G. Vegter, ``Computing Contour Generators of Evolving Implicit Surfaces,'' \emph{ACM Transactions on Graphics} \textbf{25} (2006), 1243--1280; \href{https://research.rug.nl/en/publications/computing-contour-generators-of-evolving-implicit-surfaces/}{publication page}.
\end{itemize}
\vfill\small\emph{Audit disposition: approve the corrected mathematical statements under the explicit assumptions; request originality and authorship review before journal submission.}
\end{document}
